% =============================================================================
% MIBEL — DP-02: Consumo vs. Preço
% Relatório Técnico — Data Product 02
% MEI / TEAD 2.0 | ESTG-IPP | 2025-2026
% =============================================================================
\documentclass[11pt,a4paper]{article}

\usepackage[utf8]{inputenc}
\usepackage[T1]{fontenc}
\usepackage[portuguese]{babel}
\usepackage[margin=2.2cm]{geometry}
\usepackage{graphicx}
\usepackage{booktabs}
\usepackage{longtable}
\usepackage{array}
\usepackage{xcolor}
\usepackage{listings}
\usepackage{hyperref}
\usepackage{enumitem}
\usepackage{titlesec}
\usepackage{fancyhdr}
\usepackage{microtype}
\usepackage{parskip}
\usepackage{tabularx}
\usepackage{multirow}

% Cores
\definecolor{mibel-blue}{RGB}{15,31,63}
\definecolor{mibel-accent}{RGB}{24,56,119}
\definecolor{code-bg}{RGB}{248,248,248}
\definecolor{pass-green}{RGB}{22,163,74}
\definecolor{fail-red}{RGB}{185,28,28}
\definecolor{warn-orange}{RGB}{234,88,12}

% Configuração de títulos
\titleformat{\section}{\large\bfseries\color{mibel-blue}}{}{0em}{}[\titlerule]
\titleformat{\subsection}{\normalsize\bfseries\color{mibel-accent}}{}{0em}{}
\titleformat{\subsubsection}{\normalsize\itshape}{}{0em}{}

% Cabeçalho / Rodapé
\pagestyle{fancy}
\fancyhf{}
\fancyhead[L]{\small\textcolor{mibel-blue}{\textbf{MIBEL — DP-02: Consumo vs.\ Preço}}}
\fancyhead[R]{\small\textcolor{gray}{MEI / TEAD 2.0 | ESTG-IPP}}
\fancyfoot[C]{\small\thepage}
\renewcommand{\headrulewidth}{0.4pt}

% Listings
\lstset{
  backgroundcolor=\color{code-bg},
  basicstyle=\ttfamily\footnotesize,
  breaklines=true,
  frame=single,
  rulecolor=\color{gray!30},
  xleftmargin=6pt,
  xrightmargin=6pt,
}

% Hyperref
\hypersetup{
  colorlinks=true,
  linkcolor=mibel-accent,
  urlcolor=mibel-accent,
  pdfauthor={Grupo MIBEL — MEI TEAD 2.0},
  pdftitle={MIBEL DP-02 — Relatório Técnico},
}

\setlength{\parindent}{0pt}
\setlength{\parskip}{5pt}

% =============================================================================
\begin{document}
% =============================================================================

% Capa
\begin{titlepage}
  \centering
  \vspace*{2cm}
  {\color{mibel-blue}\rule{\linewidth}{1.5pt}}\\[0.4cm]
  {\Huge\bfseries\color{mibel-blue} MIBEL\\[0.3cm]
   \Large Energy Analytics Platform}\\[0.3cm]
  {\color{mibel-blue}\rule{\linewidth}{1.5pt}}\\[1.2cm]
  {\large\textbf{Relatório Técnico — Data Product 02}}\\[0.4cm]
  {\LARGE\bfseries\color{mibel-accent} Consumo vs.\ Preço}\\[0.4cm]
  {\large\textit{Static Data Pipeline \& Streaming Data Pipeline}}\\[1cm]
  {\large Mestrado em Engenharia Informática}\\
  {\large Tecnologias e Engenharia de Dados (TEAD 2.0)}\\[0.5cm]
  {\large ESTG --- Instituto Politécnico do Porto}\\[1.5cm]
  {\large Número de Aluno: \textbf{8220800}}\\[0.3cm]
  {\large Email: \href{mailto:8220800@estg.ipp.pt}{8220800@estg.ipp.pt}}\\[2cm]
  {\large Ano Letivo 2025--2026}\\[1cm]
  \vfill
  {\small Versão 1.0 --- Maio 2026}
\end{titlepage}

\tableofcontents
\newpage

% =============================================================================
\section{Contexto e Cenário}
% =============================================================================

\subsection{Enquadramento}

O Data Product 02 (DP-02) da plataforma MIBEL responde à seguinte questão de negócio central: \textbf{qual é o custo horário da energia consumida em Portugal?} O mercado grossista ibérico (MIBEL) define um preço \textit{day-ahead} por hora, determinado diariamente na bolsa OMIE e válido para Portugal e Espanha. O preço varia substancialmente ao longo do dia e ao longo do ano — influenciado pela procura, pelo mix de produção renovável e pela interconexão ibérica — criando uma necessidade analítica de cruzar consumo real com preço de mercado numa granularidade horária.

\subsection{Cenário do Grupo}

No cenário hipotético criado, o grupo simula uma equipa de \textit{Data Engineering} de um comercializador de energia português que necessita de:

\begin{enumerate}[leftmargin=*,noitemsep]
  \item \textbf{Monitorizar} o custo estimado da energia consumida hora a hora, com base no preço OMIE;
  \item \textbf{Analisar} correlações entre padrões de consumo e variação de preço ao longo do tempo;
  \item \textbf{Prever} o consumo da hora seguinte para apoiar decisões de compra no mercado intradiário;
  \item \textbf{Receber dados em tempo quase-real} via API pública para complementar o arquivo histórico.
\end{enumerate}

Para responder a estes requisitos, o DP-02 foi implementado com \textbf{duas sub-pipelines independentes}: a \textit{Static Data Pipeline} (dados históricos OMIE em CSV desde 2023) e a \textit{Streaming Data Pipeline} (dados da API Energy-Charts / ENTSO-E em janela deslizante). As duas pipelines partilham a mesma arquitetura Medallion e o mesmo schema Gold, mas coexistem no mesmo lakehouse sem interferência.

\subsection{Fontes de Dados}

\begin{table}[h!]
\centering
\begin{tabularx}{\linewidth}{llX}
\toprule
\textbf{Pipeline} & \textbf{Fonte} & \textbf{Descrição} \\
\midrule
Static & OMIE CSV & Preços \textit{day-ahead} hora a hora (€/MWh) por dia, PT e ES \\
Static & REN/ERSE CSV & Consumo total nacional a 15 minutos (kW), BT/MT/AT/MAT \\
Streaming & Energy-Charts API & Carga horária ENTSO-E + preços OMIE via API pública (JSON) \\
\bottomrule
\end{tabularx}
\end{table}

% =============================================================================
\section{Capítulo 1: Static Data Pipeline}
% =============================================================================

\subsection{Entregáveis e Arquitetura Medallion}

A \textit{Static Data Pipeline} processa dois conjuntos de ficheiros CSV históricos e produz quatro tabelas Iceberg ao longo das camadas Bronze, Silver e Gold. A Figura~\ref{fig:static_flow} resume o fluxo.

\begin{description}[leftmargin=*]
  \item[\textbf{Bronze — Landing + Ingestão}]
    Antes da ingestão Iceberg, os CSVs são carregados para o MinIO (\texttt{s3a://warehouse/landing/}) e expostos como tabelas Hive \textit{external} (\texttt{hive.landing.consumo\_raw\_csv} e \texttt{hive.landing.preco\_raw\_csv}). Isto permite inspecção SQL dos ficheiros raw antes de qualquer transformação. Em seguida, dois \texttt{INSERT INTO} materializam os dados em Iceberg:
    \begin{itemize}[leftmargin=*,noitemsep]
      \item \texttt{iceberg.bronze.consumo\_raw} — todas as colunas da fonte (BT, MT, AT, MAT, total em kW) + \texttt{process\_date} para idempotência; particionado por \texttt{process\_date}
      \item \texttt{iceberg.bronze.preco\_raw} — preço PT e ES, hora original OMIE (1--25), \texttt{date\_raw} como string; particionado por \texttt{process\_date}
    \end{itemize}

  \item[\textbf{Silver — Normalização e Validação}]
    Dois \texttt{INSERT INTO} transformam o Bronze:
    \begin{itemize}[leftmargin=*,noitemsep]
      \item \texttt{iceberg.silver.consumo\_hourly} — agrega os registos de 15 min em 1 hora (\texttt{SUM/1000} de kW para MWh), normaliza o \texttt{datahora} para \texttt{ts\_utc}; particionado por \texttt{year/month}
      \item \texttt{iceberg.silver.preco\_hourly} — converte a hora OMIE (1--24) para \texttt{ts\_utc} via \texttt{date\_raw + (hour - 1) horas}; descarta a hora 25 (artefacto DST outono); particionado por \texttt{year/month}
    \end{itemize}

  \item[\textbf{Gold — Produto Analítico e Feature Table}]
    Dois \texttt{INSERT INTO} materializam o produto final:
    \begin{itemize}[leftmargin=*,noitemsep]
      \item \texttt{iceberg.gold.dp\_energy\_market\_hourly} — \textit{join} Silver consumo $\times$ Silver preço por \texttt{ts\_utc}; adiciona features temporais (\texttt{hora}, \texttt{dia\_semana}, \texttt{is\_weekend}), lags (\texttt{consumo\_lag\_1h}, \texttt{consumo\_lag\_24h}, \texttt{price\_lag\_1h}) e médias móveis (\texttt{rolling\_avg\_consumo\_24h}, \texttt{rolling\_avg\_price\_24h})
      \item \texttt{iceberg.gold.feat\_load\_forecasting\_hourly} — feature table ML: cópia do produto analítico com adição do \textbf{target} \texttt{consumo\_next\_hour} (LEAD de 1 hora sobre \texttt{consumo\_total})
    \end{itemize}
\end{description}

\subsection{Orquestração com Flyte (Modo Remoto)}

A \textit{Static Data Pipeline} é a única pipeline do projecto com execução \textbf{Flyte remota} (modo K3s). O orquestrador \texttt{run\_medallion\_consumo\_precos.py} auto-inicia o sandbox Flyte se não estiver em execução (build da imagem \texttt{mibel-flyte:latest} e \texttt{docker run --privileged}). Os \textit{tasks} Flyte correm em \textit{pods} K3s isolados que acedem ao Trino e ao MinIO via \texttt{host.docker.internal}.

A pipeline Flyte está dividida em módulos independentes em \texttt{workflows/}:

\begin{itemize}[leftmargin=*,noitemsep]
  \item \texttt{flyte\_bronze\_ingest.py} — upload CSV para MinIO + DDL landing + INSERT Bronze
  \item \texttt{flyte\_bronze\_to\_silver.py} — transformação Silver (agregação 15min→1h, normalização UTC)
  \item \texttt{flyte\_silver\_to\_gold.py} — joins, lags, médias móveis e feature table ML
\end{itemize}

O registo rápido (\texttt{--copy all}) carrega o código local para o bucket MinIO \texttt{flyte/}, eliminando a necessidade de publicar uma imagem Docker por cada alteração de código.

\subsection{Quality Gates — 38 Checks}

O pipeline estático é o mais instrumentado do projecto, com \textbf{38 quality checks} distribuídos pelas três camadas:

\begin{table}[h!]
\centering
\begin{tabular}{lrr}
\toprule
\textbf{Camada} & \textbf{Checks} & \textbf{Resultado típico} \\
\midrule
Bronze & 11 & 9 PASS, 2 WARN \\
Silver & 13 & 11 PASS, 2 WARN \\
Gold   & 14 & 14 PASS \\
\midrule
\textbf{Total} & \textbf{38} & \textbf{32 PASS, 6 WARN, 0 FAIL} \\
\bottomrule
\end{tabular}
\caption{Resultado típico dos quality gates na Static Data Pipeline.}
\end{table}

Os 6 WARNs observados são todos \textbf{esperados e documentados} — representam características reais do mercado energético ibérico, não erros de dados:

\begin{itemize}[leftmargin=*,noitemsep]
  \item \textbf{Preços negativos} (Bronze + Silver, 504 horas): O MIBEL permite preços \textit{day-ahead} negativos em períodos de excesso de geração renovável. O dado é válido; o WARN garante rastreabilidade. Classificado como WARN (não FAIL) porque bloquear o pipeline neste caso seria incorreto.
  \item \textbf{Duplicados em consumo\_raw} (Bronze, 24 grupos): Artefacto de re-ingestão idempotente — a segunda execução sobre o mesmo \texttt{process\_date} gera duplicados na Bronze que a Silver absorve via \texttt{GROUP BY + AVG}.
  \item \textbf{Dias de fronteira com carga parcial} (Bronze, 2 dias): Primeiros e últimos dias do ficheiro histórico podem ter menos de 96 registos de 15 min. Não afectam a agregação horária da Silver.
  \item \textbf{Dias DST com 23 horas} (Silver, 2 dias): Na transição \textit{spring forward}, o dia tem 23 horas UTC. O check \texttt{horas\_por\_dia >= 23} é WARN (não FAIL) precisamente para acomodar este comportamento correcto.
\end{itemize}

\subsection{Modelo de Machine Learning}

O modelo ML do DP-02 Static é um \textbf{GradientBoostingRegressor} (scikit-learn) que prevê \texttt{consumo\_next\_hour} (consumo da hora seguinte em MWh). As features são lidas directamente da tabela Gold \texttt{feat\_load\_forecasting\_hourly}:

\begin{itemize}[leftmargin=*,noitemsep]
  \item Features de calendário: \texttt{hora}, \texttt{dia\_semana}, \texttt{is\_weekend}
  \item Features de lag: \texttt{consumo\_lag\_1h}, \texttt{consumo\_lag\_24h}, \texttt{price\_lag\_1h}
  \item Features de janela: \texttt{rolling\_avg\_consumo\_24h}, \texttt{rolling\_avg\_price\_24h}
  \item Feature de contexto: \texttt{market\_price\_pt} (preço corrente)
\end{itemize}

O treino usa split temporal (sem \textit{shuffle}) para evitar \textit{data leakage}. Métricas, hiperparâmetros e artefactos (\texttt{feature\_importances.json}, modelo sklearn serializado) são registados no MLflow com tags de rastreabilidade (\texttt{dataset=iceberg.gold.feat\_load\_forecasting\_hourly}, \texttt{target=consumo\_next\_hour}).

\subsection{Dashboard Grafana — Static}

O dashboard \texttt{consumo\_preco\_overview.json} é provisionado automaticamente e expõe:

\begin{itemize}[leftmargin=*,noitemsep]
  \item KPI cards: consumo médio horário (MWh), preço médio (€/MWh), custo estimado (€)
  \item Série temporal dual-axis: consumo e preço por hora
  \item Distribuição mensal de preços (boxplot/heatmap)
  \item Análise de correlação preço--consumo por período do dia e dia de semana
  \item Painel de previsão: consumo previsto vs.\ real (via inferência do modelo MLflow)
\end{itemize}

O datasource Grafana executa SQL direto sobre Trino — sem ETL adicional — usando a tabela Gold como única fonte de verdade.

\subsection{Decisões de Impacto na Qualidade e Performance}

\subsubsection{Hora 25 OMIE — Descarte Controlado na Silver}

Os ficheiros OMIE incluem ocasionalmente uma ``hora 25'' nos dias de transição DST de outono (relógio recua 1 hora). Preservar esta hora como-está introduziria um timestamp ambíguo no Silver. A decisão foi \textbf{descartar a hora 25 na Silver} via filtro \texttt{WHERE hour <= 24}. O impacto: perda de 1 ponto de dados por ano de transição DST — aceitável face ao risco de ambiguidade UTC.

\subsubsection{Preços Negativos — WARN, não FAIL}

Classificar preços negativos como FAIL bloquearia o pipeline em situações de mercado legítimas (excesso renovável). A decisão de WARN permite que o dado chegue ao Gold e ao dashboard, onde é visível e rastreável, sem impedir o fluxo de dados. Esta distinção entre anomalia de mercado e erro de dados é uma decisão de engenharia deliberada.

\subsubsection{Feature Tables Pré-computadas no Gold}

Os lags e médias móveis são calculados uma única vez em SQL Trino e materializados em Iceberg na tabela \texttt{feat\_load\_forecasting\_hourly}. A alternativa (calcular em Python durante o treino) introduziria risco de \textit{training-serving skew}: qualquer diferença mínima entre o cálculo em treino e o cálculo em produção degrada as previsões silenciosamente. A pré-computação no Gold garante que o modelo em treino e o modelo em serving usam exactamente a mesma lógica SQL.

\subsubsection{Landing Zone como Camada de Observabilidade}

A decisão de expor os CSVs raw como tabelas Hive \textit{external} antes da ingestão Iceberg tem um impacto de qualidade: qualquer analista pode fazer \texttt{SELECT} sobre \texttt{hive.landing.consumo\_raw\_csv} ou \texttt{hive.landing.preco\_raw\_csv} e comparar com o Bronze para auditar a transformação. Isto não adiciona custo de armazenamento (os ficheiros já existem) mas oferece observabilidade completa do processo de ingestão.

\subsection{Relevância para o Negócio}

\subsubsection{Custo Horário como Métrica Operacional}

A coluna derivada \texttt{consumo\_total × market\_price\_pt} na Gold corresponde ao \textbf{custo estimado da energia consumida em Portugal em cada hora}. Para um comercializador, esta métrica é o KPI primário da margem bruta: a diferença entre o preço de venda ao cliente final e o preço de compra no mercado OMIE determina a rentabilidade hora a hora.

\subsubsection{Previsão de Consumo para Optimização de Portfolio}

O modelo GBR com target \texttt{consumo\_next\_hour} oferece um sinal de decisão directo: se o consumo previsto para a próxima hora é significativamente superior ao actual, o comercializador pode comprar energia antecipadamente no mercado intradiário (tipicamente mais barato do que o mercado de balanço de curto prazo). O lag de 24 horas (\texttt{consumo\_lag\_24h}) é tipicamente a feature mais importante — captura o padrão semanal de consumo — o que reflecte a forte sazonalidade diária/semanal do consumo eléctrico português.

\subsubsection{Correlação Preço--Consumo como Sinal de Hedging}

A presença de \texttt{price\_lag\_1h} e \texttt{rolling\_avg\_price\_24h} nas features permite ao modelo capturar a relação entre preço e consumo: em horas de preço muito elevado, o consumo industrial tende a reduzir (resposta à procura). Analisar esta correlação no dashboard ajuda o departamento comercial a identificar clientes com elasticidade ao preço e a desenhar produtos com tarifas dinâmicas.

% =============================================================================
\section{Capítulo 2: Streaming Data Pipeline}
% =============================================================================

\subsection{Entregáveis e Arquitetura Medallion}

A \textit{Streaming Data Pipeline} obtém dados em tempo quase-real via \textbf{Energy-Charts API} (wrapper público da ENTSO-E Transparency Platform) e executa a mesma arquitectura Medallion do pipeline estático — mas com tabelas com sufixo \texttt{\_api} para coexistência no mesmo lakehouse. O orquestrador \texttt{run\_streaming\_pipeline.py} suporta janelas flexíveis (\texttt{--days N}, \texttt{--start/--end}, \texttt{--today}, \texttt{--full}).

\begin{description}[leftmargin=*]
  \item[\textbf{Bronze — Ingestão API}]
    O script Python faz chamadas à Energy-Charts API e insere directamente em Iceberg (sem landing zone CSV):
    \begin{itemize}[leftmargin=*,noitemsep]
      \item \texttt{iceberg.bronze.consumo\_api\_raw} — carga horária nacional em MW (\texttt{total}), \texttt{ts\_utc} já normalizado na ingestão, \texttt{source\_url} para rastreabilidade, \texttt{fetch\_date} e \texttt{process\_date}
      \item \texttt{iceberg.bronze.preco\_api\_raw} — preço PT e ES em €/MWh, \texttt{ts\_utc} normalizado; \texttt{price\_spain\_eur\_mwh} pode ser NULL (endpoint não garante ES em todos os períodos)
    \end{itemize}
    \textbf{Diferença face ao pipeline estático:} o timestamp UTC já está normalizado na ingestão (derivado de \texttt{unix\_seconds} da resposta JSON), eliminando a necessidade de parsing de strings na Silver.

  \item[\textbf{Silver — Validação e Deduplicação}]
    \begin{itemize}[leftmargin=*,noitemsep]
      \item \texttt{iceberg.silver.consumo\_api\_hourly} — filtra nulos, converte MW → MWh (a fonte já é horária, a conversão é identidade: \texttt{total × 1.0}); particionado por \texttt{year/month}
      \item \texttt{iceberg.silver.preco\_api\_hourly} — deduplicação por \texttt{ts\_utc} via \texttt{GROUP BY + AVG} (a API pode retornar sobreposição em re-ingestões); particionado por \texttt{year/month}
    \end{itemize}

  \item[\textbf{Gold — Produto Analítico e Feature Table}]
    Schema \textbf{idêntico} ao pipeline estático, com sufixo \texttt{\_api}:
    \begin{itemize}[leftmargin=*,noitemsep]
      \item \texttt{iceberg.gold.dp\_energy\_market\_api\_hourly} — join consumo × preço + features temporais + lags + médias móveis
      \item \texttt{iceberg.gold.feat\_load\_forecasting\_api\_hourly} — feature table ML com \texttt{consumo\_next\_hour} (LEAD 1h)
    \end{itemize}
\end{description}

\subsection{Orquestração Local (sem Flyte)}

Ao contrário do pipeline estático, a \textit{Streaming Data Pipeline} corre em modo local: o orquestrador Python executa as \textit{tasks} directamente no processo local, sem Flyte Sandbox. Esta decisão é motivada pela natureza da operação: a API é consultada regularmente (ex.: diariamente) com volumes de dados pequenos (dias ou semanas), não justificando o overhead do K3s. O Flyte é reservado para o pipeline em batch sobre grandes volumes históricos.

O orquestrador suporta flags granulares compatíveis com o pipeline estático:

\begin{lstlisting}
python run_streaming_pipeline.py --days 7 --skip-docker
python run_streaming_pipeline.py --start 2024-01-01 --end 2024-12-31 --skip-docker
python run_streaming_pipeline.py --today --skip-docker --no-quality
\end{lstlisting}

\subsection{Pré-requisito: Token ENTSO-E}

A API Energy-Charts requer autenticação via \textbf{ENTSOE\_TOKEN} (token gratuito pedido por email a \texttt{transparency@entsoe.eu}). O orquestrador verifica a presença da variável de ambiente antes de iniciar e falha imediatamente com mensagem clara se o token estiver ausente.

\begin{lstlisting}
# PowerShell
$env:ENTSOE_TOKEN = "<token>"
python run_streaming_pipeline.py --days 7 --skip-docker
\end{lstlisting}

\subsection{Quality Gates — Streaming}

O pipeline API tem quality checks adaptados à sua fonte dinâmica:

\begin{table}[h!]
\centering
\begin{tabular}{lrp{7cm}}
\toprule
\textbf{Camada} & \textbf{Checks} & \textbf{Foco principal} \\
\midrule
Bronze & $\sim$9 & Nulos em \texttt{ts\_utc/total/price}, range de preços, unicidade por \texttt{ts\_utc} \\
Silver & $\sim$10 & Completude horária ($\geq$23h/dia), deduplicação, range MWh \\
Gold   & $\sim$11 & Integridade do join, lags não-nulos, freshness (últimas 24h presentes) \\
\bottomrule
\end{tabular}
\end{table}

Um check específico do pipeline API é o de \textbf{freshness}: verifica se existem dados das últimas 24 horas no Gold, garantindo que a janela deslizante está actualizada. No pipeline estático, este check não faz sentido (dados históricos fixos).

\subsection{Dashboard Grafana — Streaming}

O dashboard \texttt{consumo\_preco\_streaming\_overview.json} expõe dados da fonte API em tempo quase-real:

\begin{itemize}[leftmargin=*,noitemsep]
  \item KPI cards com dados das últimas 24 horas: consumo atual, preço spot corrente, tendência de consumo vs.\ hora anterior
  \item Série temporal das últimas 72 horas (consumo + preço dual-axis)
  \item Painel de comparação entre fontes: sobreposição de \texttt{dp\_energy\_market\_hourly} (static) e \texttt{dp\_energy\_market\_api\_hourly} (streaming) para os períodos em comum
  \item Alerta visual quando \texttt{market\_price\_pt} supera limiar configurável (ex.: 100 €/MWh)
\end{itemize}

\subsection{Decisões de Impacto na Qualidade e Performance}

\subsubsection{Coexistência via Sufixo \texttt{\_api}}

A decisão de usar tabelas com sufixo \texttt{\_api} (\texttt{consumo\_api\_raw}, \texttt{dp\_energy\_market\_api\_hourly}, etc.) em vez de escrever nas mesmas tabelas do pipeline estático tem impacto directo na qualidade e na operação:

\begin{itemize}[leftmargin=*,noitemsep]
  \item \textbf{Isolamento de falhas}: uma API indisponível ou com dados incorrectos não afecta os dados históricos CSV já validados
  \item \textbf{Comparabilidade}: o mesmo período temporal pode ser consultado em ambas as fontes para validar coerência (PT vs.\ ENTSO-E vs.\ OMIE)
  \item \textbf{Schema idêntico no Gold}: o dashboard e os modelos ML podem alternar entre fontes sem modificação — apenas a referência à tabela muda
\end{itemize}

\subsubsection{Normalização UTC na Ingestão (Bronze)}

No pipeline estático, o timestamp UTC é derivado na Silver (a partir de \texttt{date\_raw} + \texttt{hour}). No pipeline API, o \texttt{ts\_utc} é calculado na ingestão Bronze (a partir do \texttt{unix\_seconds} da resposta JSON). Esta decisão encurta o fluxo de transformação e elimina o risco de erro de parsing de data/hora: a API devolve \texttt{unix\_seconds} não-ambíguo, enquanto os CSVs OMIE usam hora local com DST.

\subsubsection{Deduplicação Silver via GROUP BY + AVG}

A API Energy-Charts pode retornar horas sobrepostas em re-ingestões dentro da mesma janela (ex.: executar \texttt{--days 7} dois dias seguidos gera 6 dias de sobreposição). A Silver resolve isto via \texttt{GROUP BY ts\_utc + AVG(valor)}, absorvendo duplicados sem propagação para o Gold. O quality check Bronze \textit{uniqueness} emite WARN (não FAIL) para documentar a ocorrência.

\subsubsection{Sem Flyte no Streaming — Decisão de Proporcionalidade}

Manter a Streaming Pipeline em modo local (sem K3s) foi uma decisão deliberada de proporcionalidade técnica: o overhead de provisionar e gerir um sandbox Flyte para operações de minutos com poucos milhares de linhas não tem justificação operacional. O Flyte justifica-se no pipeline estático, onde o volume é substancialmente maior (anos de dados históricos, centenas de milhares de linhas) e a orquestração distribuída traz benefícios reais de paralelismo e gestão de falhas por \textit{task}.

\subsection{Relevância para o Negócio}

\subsubsection{Dados em Tempo Quase-Real para Decisões Intradiárias}

A ENTSO-E disponibiliza os dados de carga actual com latência de 15--60 minutos. Um comercializador que execute \texttt{run\_streaming\_pipeline.py --today} a cada hora obtém um Gold actualizado que alimenta automaticamente o dashboard Grafana. Este ciclo de actualização suporta decisões intradiárias: o operador de mercado pode visualizar se o consumo actual está acima ou abaixo da média histórica e ajustar a posição de compra/venda.

\subsubsection{Validação Cruzada entre Fontes como Indicador de Qualidade}

O painel de comparação entre fontes no dashboard Streaming tem um valor de negócio específico: permite detectar divergências entre os dados OMIE CSV (fonte oficial do mercado) e os dados Energy-Charts/ENTSO-E (fonte API). Uma divergência persistente sinaliza um problema de fonte (ex.: atraso de publicação na API, revisão retroactiva dos dados OMIE) que pode afectar a fiabilidade das previsões.

\subsubsection{Extensibilidade para Alertas Automáticos}

A arquitectura Streaming do DP-02 é a base para um sistema de alertas automáticos hipotético: quando \texttt{market\_price\_pt} supera um limiar crítico (ex.: 200 €/MWh, valor observado em crises de gás natural), o pipeline pode desencadear uma notificação ao departamento comercial para activar contratos de hedging pré-acordados. O Redpanda (broker Kafka-compatível) já presente na stack Docker é o componente natural para este fluxo de eventos.

% =============================================================================
\section{Síntese Comparativa das Duas Sub-Pipelines}
% =============================================================================

\begin{table}[h!]
\centering
\small
\begin{tabularx}{\linewidth}{lXX}
\toprule
\textbf{Dimensão} & \textbf{Static Data Pipeline} & \textbf{Streaming Data Pipeline} \\
\midrule
Fonte & CSVs OMIE + REN/ERSE & Energy-Charts API (ENTSO-E) \\
Cobertura temporal & Histórico completo (2023--) & Janela deslizante (\texttt{--days N}) \\
Orquestração & Flyte remoto (K3s) & Local (sem Flyte) \\
Normalização UTC & Silver (parsing date\_raw + hour) & Bronze (unix\_seconds da API) \\
DST handling & Descarte hora 25 OMIE & Não aplicável (API devolve UTC) \\
Preços negativos & WARN (504 horas observadas) & WARN (idem) \\
Quality checks & 38 (11+13+14) & $\sim$30 (9+10+11) \\
Tabelas Gold & \texttt{dp\_energy\_market\_hourly} & \texttt{dp\_energy\_market\_api\_hourly} \\
Feature table ML & \texttt{feat\_load\_forecasting\_hourly} & \texttt{feat\_load\_forecasting\_api\_hourly} \\
Dashboard & \texttt{consumo\_preco\_overview} & \texttt{consumo\_preco\_streaming\_overview} \\
\bottomrule
\end{tabularx}
\caption{Comparação entre as duas sub-pipelines do DP-02.}
\end{table}

A principal decisão arquitectural do DP-02 é a \textbf{coexistência das duas pipelines no mesmo lakehouse com schema Gold idêntico}. Esta escolha garante que o dashboard, os modelos ML e as APIs consomem dados de qualquer das fontes sem modificação --- apenas a referência à tabela muda. Em ambiente de produção, a Static Pipeline fornece o arquivo histórico de alta fidelidade (dados OMIE oficiais) enquanto a Streaming Pipeline fornece a janela operacional de curto prazo com menor latência.

% =============================================================================
\section{Conclusão}
% =============================================================================

O Data Product 02 demonstra que uma arquitectura Medallion pode acomodar duas fontes heterogéneas --- batch histórico em CSV e API em tempo quase-real --- no mesmo lakehouse, com quality gates diferenciados por camada e fonte, sem comprometer a coerência do schema de consumo downstream.

As decisões de maior impacto foram:

\begin{itemize}[leftmargin=*,noitemsep]
  \item \textbf{Preços negativos como WARN} — distingue anomalia de mercado de erro de dados, evitando bloqueios indevidos do pipeline
  \item \textbf{Descarte controlado da hora 25 OMIE na Silver} — elimina ambiguidade UTC sem perda de qualidade analítica
  \item \textbf{Feature tables pré-computadas no Gold} — elimina \textit{training-serving skew} e garante reprodutibilidade dos modelos ML
  \item \textbf{Sufixo \texttt{\_api} para coexistência} — isola falhas de fonte e permite comparação entre CSV oficial e API em tempo real
  \item \textbf{Flyte remoto apenas no Static Pipeline} — proporcionalidade técnica: orquestração distribuída onde o volume o justifica
\end{itemize}

Do ponto de vista de negócio, o DP-02 oferece ao comercializador de energia a capacidade de monitorizar o custo horário da energia em tempo quase-real, prever o consumo da hora seguinte para optimização de portfolio e detectar automaticamente condições de mercado extremas (preços negativos, picos de preço) que requerem acção imediata.

% =============================================================================
\end{document}
